%-------------------------------------------------------------------------------
%	ABSTRACT PAGE
%-------------------------------------------------------------------------------

\addchaptertocentry{Хураангуй} % Хураангуйг гарчигт нэмэх

\begin{center}
{\scshape\Large \univname\par} % Их сургуулийн нэр
{\scshape\large \facname\par}\vspace{0.5cm} % Их сургуулийн нэр
{\huge\textbf{{Хураангуй}} \par}
\bigskip
{\Large{\ttitle} \par} % Тезисийн нэр
\bigskip

{\normalsize \shortname \par} % Зохиогчийн нэр
\addressname
\end{center}

\textit{\textbf{Түлхүүр үгс: \keywordnames}}
\bigskip

Монголын загварын салбар, дизайнерууд, улирлын коллекцууд, редакцийн нийтлэлүүд болон зурагт материалууд нь цахим орчинд ихэвчлэн салангид байрлаж, хайх, ангилах, засварлах, дахин ашиглахад хүндрэлтэй байдаг. Загварын брэндүүдийн танилцуулга, коллекцын look бүрийн материал, улирал, нийтлэл, ярилцлага, архивын зураг зэрэг мэдээллийг нэг загвартай өгөгдлийн платформд хадгалах нь соёлын өв, бүтээлч үйлдвэрлэл, редакцийн ажлын хувьд чухал шаардлага болж байна.

Энэхүү төгсөлтийн ажлын хүрээнд Монголын загварын архивт зориулсан Anoce нэртэй веб платформыг боловсруулж, хөгжүүлэв. Уг платформ нь дизайнер ба брэндийн лавлах, улирлын коллекц, look зураг, редакцийн нийтлэл, глобал хайлт, хэрэглэгчийн bookmark/follow боломж, админ контент удирдлага, asset upload, хэрэглэгчийн эрхийн хяналт зэрэг үндсэн модулиудаас бүрдэнэ. Frontend нь Next.js App Router, React, TypeScript, Tailwind CSS болон Framer Motion дээр суурилж, public archive болон editorial туршлагыг responsive, визуал төвтэй байдлаар хүргэнэ. Контентын үндсэн repository нь CouchDB дээр document хэлбэрээр хадгалагдаж, Supabase нь authentication, profile role, bookmark болон RAG индексийн нэмэлт үүргийг гүйцэтгэнэ.

Хиймэл оюун ухааны хэсэг нь архивын контентод тулгуурласан туслах функц байдлаар хэрэгжсэн. \texttt{/api/chat} route нь live archive search, Supabase RAG document search, prompt бүрдүүлэлт, Gemini/Ollama provider fallback ашиглан хэрэглэгчийн загвар, дизайнер, collection, trend-тэй холбоотой асуултад контексттэй хариулт өгөх боломжтой. Иймээс уг ажил нь Монголын загварын мэдээллийг зөвхөн танилцуулах вебсайт бус, харин контент удирдлага, хайлт, архивын өгөгдөл, AI туслах, хэрэглэгчийн хувийн үйлдлийг нэгтгэсэн дижитал архивын платформ болж байгаагаараа практик ач холбогдолтой.
